\documentclass[12pt,a4paper]{article}
\usepackage[utf8]{inputenc}
\usepackage[T1]{fontenc}
\usepackage[spanish]{babel}
\usepackage[margin=1.4cm]{geometry}
\usepackage{array,longtable}
\usepackage{titlesec}
\usepackage{xcolor}
\usepackage{hyperref}

\definecolor{HeadingColor}{HTML}{243B53}
\definecolor{InstructionColor}{HTML}{704214}
\definecolor{ConceptColor}{HTML}{205C50}
\definecolor{LinkColor}{HTML}{0047AB}
\hypersetup{colorlinks=true,allcolors=LinkColor,pdfborder={0 0 0}}
% Los textos entre corchetes son campos editables.
\newcommand{\campo}[1]{\textcolor{InstructionColor}{\textit{[#1]}}}
\newcommand{\concept}[1]{\textcolor{ConceptColor}{\textbf{#1}}}
\newcommand{\respuesta}[2]{\par\textbf{#1}\par\campo{#2}\par}
\setlength{\parindent}{0pt}
\setlength{\parskip}{3pt}
\setlength{\tabcolsep}{4pt}
\renewcommand{\arraystretch}{1.15}
\setlength{\LTpre}{4pt}
\setlength{\LTpost}{4pt}
\titleformat{\section}{\normalsize\bfseries\color{HeadingColor}}{}{0pt}{}
\titlespacing*{\section}{0pt}{8pt}{3pt}
\newcolumntype{P}[1]{>{\raggedright\arraybackslash}p{#1}}
% El contador conserva D1 como referencia al incorporar estas secciones al proyecto.
\newcounter{desafioCsiete}
\renewcommand{\thedesafioCsiete}{D\arabic{desafioCsiete}}

\begin{document}
\raggedbottom
\begin{center}
{\Large\bfseries\color{HeadingColor} C7--C8: Definición del desafío y conexiones disciplinares}\\[3pt]
Emprendimiento e Innovación --- Entrega del proyecto
\end{center}
\respuesta{Proyecto}{Nombre del proyecto}
\respuesta{Integrantes}{Nombres de los estudiantes}
\textbf{Grado:} \campo{Grado}\hfill\textbf{Fecha de actualización:} \campo{Fecha}

\section*{C7 --- 1. Contexto del desafío}
{\small
\begin{longtable}{|P{0.23\linewidth}|P{\dimexpr0.77\linewidth-4\tabcolsep-3\arrayrulewidth\relax}|}
\hline
\textbf{Elemento} & \textbf{Información del proyecto} \\\hline
\endfirsthead
\hline\textbf{Elemento} & \textbf{Información del proyecto (continuación)} \\\hline
\endhead
Oportunidad priorizada & \campo{Código de oportunidad: A, B o C}; \campo{Descripción breve de la oportunidad seleccionada en C6} \\\hline
Hallazgos relacionados & \campo{Hallazgos relacionados: H1, H2 u otros; referencia a su registro en C5} \\\hline
Estado de la selección & \campo{Definitiva o provisional; dato pendiente y efecto sobre la delimitación, si corresponde} \\\hline
Usuarios específicos & \campo{Usuarios y características pertinentes del grupo que experimenta la situación} \\\hline
Necesidad o dificultad & \campo{Necesidad: qué necesita lograr el grupo y qué dificultad se lo impide} \\\hline
Contexto y alcance & \campo{Contexto: lugar, momentos y condiciones; parte de la situación abordada y límites de lo que queda fuera} \\\hline
Evidencia de respaldo & \campo{Observación o respuesta registrada; códigos H, fuente, actividad, registro y fecha disponible; vínculo con usuarios, necesidad y contexto} \\\hline
\end{longtable}
}
\respuesta{Delimitación pendiente}{Supuesto o dato todavía sin respaldo y su efecto sobre el alcance; no aplica si corresponde}

\section*{C7 --- 2. Formulación inicial}
\textbf{Desafío inicial:} ¿Cómo podríamos \campo{mejora deseada} para \campo{usuarios específicos} en \campo{contexto delimitado}?

\section*{C7 --- 3. Revisión y ajustes}
\respuesta{Formulación inicial revisada}{Referencia a la formulación inicial anterior y parte de la pregunta examinada}
\respuesta{Comprensión compartida}{Interpretación acordada de quién necesita qué y en qué contexto; diferencias de interpretación identificadas}
\respuesta{Ambigüedad, supuesto o limitación identificada}{Término ambiguo, supuesto sin evidencia, límite de alcance o restricción anticipada a una solución; hallazgo o registro relacionado; ninguna si corresponde}
\respuesta{Ajuste realizado y justificación}{Cambio respecto a la formulación inicial y razón sustentada; si no hubo cambios, fundamento para conservarla}
\respuesta{Criterios de la formulación acordada}{Razones por las que la pregunta es clara, específica, comprensible, respaldada por los hallazgos y abierta a distintas respuestas sin imponer un producto o tecnología}

\section*{C7 --- 4. Desafío final acordado}
\refstepcounter{desafioCsiete}\label{desafio:C7:D1}
\concept{\thedesafioCsiete:} ¿Cómo podríamos \campo{mejora deseada acordada} para \campo{usuarios específicos acordados} en \campo{contexto delimitado acordado}?

\section*{C8 --- 1. Conexiones interdisciplinarias}
\textbf{Desafío de referencia:} \ref{desafio:C7:D1}.
{\small
\begin{longtable}{|P{0.09\linewidth}|P{0.22\linewidth}|P{0.22\linewidth}|P{0.25\linewidth}|P{\dimexpr0.22\linewidth-10\tabcolsep-6\arrayrulewidth\relax}|}
\hline
\textbf{Código} & \textbf{Disciplina y conocimiento específico} & \textbf{Aporte a la comprensión de D1} & \textbf{Justificación contextual} & \textbf{Fuente consultada} \\\hline
\endfirsthead
\hline
\textbf{Código} & \textbf{Disciplina y conocimiento específico} & \textbf{Aporte a la comprensión de D1} & \textbf{Justificación contextual} & \textbf{Fuente consultada} \\\hline
\endhead
\concept{K1} & \campo{Disciplina 1 y concepto concreto establecido} & \campo{Qué permite comprender o investigar sobre D1} & \campo{Hallazgo H o condición registrada que hace pertinente el concepto; fuente contextual} & \campo{Autor o entidad, título, fecha y enlace o referencia localizable} \\\hline
\concept{K2} & \campo{Disciplina 2 y concepto concreto establecido} & \campo{Qué permite comprender o investigar sobre D1} & \campo{Hallazgo H o condición registrada que hace pertinente el concepto; fuente contextual} & \campo{Autor o entidad, título, fecha y enlace o referencia localizable} \\\hline
\end{longtable}
}
% Estos párrafos admiten respuestas extensas y pueden continuar en otra página.
\respuesta{K1 --- Conocimiento consultado y alcance}{Proposición o criterio respaldado por la fuente disciplinar; alcance y límites de su aplicación a D1}
\respuesta{K1 --- Supuesto sobre el contexto}{Hipótesis contextual todavía sin demostrar, distinta del conocimiento consultado; evidencia que falta para verificarla; ninguno si corresponde}
\respuesta{K2 --- Conocimiento consultado y alcance}{Proposición o criterio respaldado por la fuente disciplinar; alcance y límites de su aplicación a D1}
\respuesta{K2 --- Supuesto sobre el contexto}{Hipótesis contextual todavía sin demostrar, distinta del conocimiento consultado; evidencia que falta para verificarla; ninguno si corresponde}

\section*{C8 --- 2. Complementariedad y conocimiento pendiente}
\respuesta{Complementariedad de K1 y K2}{Relación entre ambos aportes y comprensión conjunta de D1 que permite esa relación}
\respuesta{Conocimiento o información pendiente}{Pregunta por investigar o verificar; vínculo con D1, K1 o K2 y hallazgos H; distinción entre conocimiento disciplinar faltante y supuesto contextual pendiente}
\respuesta{Respaldo necesario}{Información, registro o fuente por consultar para resolver la pregunta pendiente y criterio para contrastar el supuesto}

\end{document}


